\chapter{فرم‌های پلکانی و حل دستگاه‌های خطی}
\label{ch:echelon}

در فصل قبل دیدیم که عملیات سطری جواب‌های یک دستگاه را تغییر نمی‌دهند. حالا سؤال این است: \emph{با این عملیات، ماتریس را به چه شکلی برسانیم تا جواب‌ها «خودشان را نشان دهند»؟} پاسخ، فرم‌های پلکانی است. در این فصل الگوریتم کامل حل دستگاه (حذف گاوسی)، شرط وجود و یکتایی جواب، تجزیه‌ی $LU$ و تعبیر هندسی مجموعه‌ی جواب را می‌بینیم.

\section{ماتریس سطری‌-تحویل‌یافته}

\begin{definition}[ماتریس سطری‌-تحویل‌یافته]
یک ماتریس $m\times n$ مثل $R$ را \term{سطری‌-تحویل‌یافته}{Row-Reduced} می‌نامیم اگر:
\begin{enumerate}
  \item \textbf{درایه‌های پیشرو برابر ۱ باشند:} اولین درایه‌ی ناصفر هر سطر ناصفرِ $R$ برابر ۱ باشد. (به اولین درایه‌ی ناصفر یک سطر، \term{درایه‌ی پیشرو}{Leading Entry} آن سطر می‌گوییم.)
  \item \textbf{بقیه‌ی درایه‌های ستون‌های پیشرو صفر باشند:} هر ستونی از $R$ که درایه‌ی پیشروِ یک سطر را در خود دارد، همه‌ی درایه‌های دیگرش صفر باشد.
\end{enumerate}
\end{definition}

\begin{example}[سطری‌-تحویل‌یافته هست یا نه؟]
کدام ماتریس‌ها سطری‌-تحویل‌یافته‌اند؟
\[
A=I_n,\qquad
B=\begin{bmatrix}1&0&0\\0&2&-1\\0&0&0\end{bmatrix},\qquad
C=\begin{bmatrix}1&0&-3\\0&1&0\\0&0&1\end{bmatrix},\qquad
D=\begin{bmatrix}1&0&0\\0&0&0\\0&0&-6\end{bmatrix}.
\]
\end{example}
\begin{solution}
\begin{itemize}
  \item $A=I_n$: \textcolor{thmcolor}{\textbf{بله.}} درایه‌ی پیشروی هر سطر ۱ است و تنها درایه‌ی ناصفر ستون خودش است.
  \item $B$: \textcolor{warncolor}{\textbf{خیر.}} درایه‌ی پیشروی سطر دوم ۲ است، نه ۱. (اصلاح: سطر ۲ را بر ۲ تقسیم کنید.)
  \item $C$: \textcolor{warncolor}{\textbf{خیر.}} ستون سوم درایه‌ی پیشروی سطر سوم را دارد، اما در سطر اول آن ستون $-3$ نشسته است. (اصلاح: $R_1\leftarrow R_1+3R_3$.)
  \item $D$: \textcolor{warncolor}{\textbf{خیر.}} درایه‌ی پیشروی سطر سوم $-6$ است. (اصلاح: سطر ۳ را بر $-6$ تقسیم کنید.) دقت کنید که وجود سطر صفر \emph{وسط} ماتریس در این تعریف مشکلی ایجاد نمی‌کند؛ این شرط در تعریف فرم پلکانی می‌آید.
\end{itemize}
\end{solution}

\begin{theorem}\label{thm:row-reduced-exists}
هر ماتریس $m\times n$ با یک ماتریس سطری‌-تحویل‌یافته هم‌ارز سطری است.
\end{theorem}
\begin{proof}
اثبات در واقع یک الگوریتم است: سطرها را یکی‌یکی از بالا به پایین «تمیز» می‌کنیم.

\textbf{سطر اول.} اگر سطر اول تمام‌صفر است، کاری نمی‌کنیم. وگرنه فرض کنید $k$ اولین ستونی باشد که $A_{1k}\neq0$. سطر اول را در $A_{1k}^{-1}$ ضرب می‌کنیم تا درایه‌ی پیشرو ۱ شود. سپس برای هر $i\neq1$، عمل $R_i\leftarrow R_i-A_{ik}R_1$ را انجام می‌دهیم. حالا ستون $k$ فقط در سطر اول ناصفر (و برابر ۱) است.

\textbf{گام کلی.} فرض کنید سطرهای $1,\dots,r$ پردازش شده‌اند؛ یعنی هر سطر ناصفر از میان آن‌ها درایه‌ی پیشروی ۱ دارد و ستون آن درایه‌ی پیشرو، جای دیگری ناصفر نیست. سطر $r+1$ را در نظر بگیرید. اگر صفر است، سراغ سطر بعد می‌رویم. وگرنه فرض کنید درایه‌ی پیشروی آن در ستون $k'$ باشد. سطر $r+1$ را بر این درایه تقسیم می‌کنیم و سپس با عمل‌های جایگزینی، درایه‌های ستون $k'$ را در \emph{همه‌ی} سطرهای دیگر صفر می‌کنیم.

\textbf{چرا کار قبلی خراب نمی‌شود؟} دو چیز را باید بررسی کنیم:
\begin{itemize}
  \item سطر $r+1$ در ستون‌های پیشروِ سطرهای $1,\dots,r$ صفر است (چون آن ستون‌ها قبلاً تمیز شده‌اند). پس اضافه کردن مضربی از سطر $r+1$ به سطرهای دیگر، درایه‌های آن ستون‌ها را تغییر نمی‌دهد.
  \item فرض کنید سطر $i\le r$ درایه‌ی پیشرویی در ستون $k_i$ دارد. اگر $k'<k_i$ باشد، درایه‌ی سطر $i$ در ستون $k'$ صفر است (چون قبل از درایه‌ی پیشروی آن است)؛ پس ضریبی که با آن سطر $r+1$ را از سطر $i$ کم می‌کنیم صفر است و سطر $i$ اصلاً تغییر نمی‌کند. اگر $k'>k_i$ باشد، سطر $r+1$ فقط در ستون‌های $\ge k'$ ناصفر است، پس فقط درایه‌های سطر $i$ در ستون‌های بعد از $k_i$ عوض می‌شوند و درایه‌ی پیشروی سطر $i$ سر جایش می‌ماند.
\end{itemize}
پس بعد از پردازش سطر $r+1$، سطرهای $1,\dots,r+1$ دو شرط تعریف را دارند. پس از $m$ گام، ماتریس سطری‌-تحویل‌یافته است. چون فقط از عملیات سطری مقدماتی استفاده کردیم، نتیجه با $A$ هم‌ارز سطری است.
\end{proof}

\section{فرم پلکانی}

\begin{definition}[فرم پلکانی]
یک ماتریس مستطیلی در \term{فرم پلکانی}{Echelon Form / Row Echelon Form} است اگر سه ویژگی زیر را داشته باشد:
\begin{enumerate}
  \item همه‌ی سطرهای ناصفر بالای همه‌ی سطرهای تمام‌صفر باشند.
  \item درایه‌ی پیشروی هر سطر، در ستونی \emph{سمت راستِ} ستون درایه‌ی پیشروی سطر بالایی باشد.
  \item همه‌ی درایه‌های زیر یک درایه‌ی پیشرو، صفر باشند.
\end{enumerate}
\end{definition}

مثلاً ماتریس زیر در فرم پلکانی است (درایه‌های پیشرو قاب گرفته شده‌اند):
\[
\begin{bmatrix}
\boxed{2}&-3&2&1\\
0&\boxed{1}&-4&8\\
0&0&0&\boxed{5/2}
\end{bmatrix}.
\]
اسم «پلکانی» از همین شکل می‌آید: درایه‌های پیشرو مثل پله‌هایی هستند که به سمت راست و پایین می‌روند.

\begin{center}
\begin{tikzpicture}[scale=0.55]
  \fill[primary!10] (0,4)--(0,3)--(2,3)--(2,2)--(3,2)--(3,1)--(6,1)--(6,0)--(7,0)--(7,4)--cycle;
  \draw[thick, primary] (0,3)--(2,3)--(2,2)--(3,2)--(3,1)--(6,1)--(6,0)--(7,0);
  \foreach \x/\y in {0.5/3.5, 2.5/2.5, 3.5/1.5, 6.5/0.5} \node[fill=excolor, text=white, inner sep=1.5pt, font=\scriptsize] at (\x,\y) {$\blacksquare$};
  \foreach \x/\y in {1.5/3.5, 2.5/3.5, 3.5/3.5, 4.5/3.5, 5.5/3.5, 6.5/3.5, 3.5/2.5, 4.5/2.5, 5.5/2.5, 6.5/2.5, 4.5/1.5, 5.5/1.5, 6.5/1.5} \node[font=\small] at (\x,\y) {$*$};
  \foreach \x/\y in {0.5/2.5,0.5/1.5,1.5/1.5,2.5/1.5,0.5/0.5,1.5/0.5,2.5/0.5,3.5/0.5,4.5/0.5,5.5/0.5,1.5/2.5} \node[gray, font=\small] at (\x,\y) {$0$};
  \draw[gray] (0,0) rectangle (7,4);
\end{tikzpicture}
\\[1mm]{\small \fa{مربع‌های نارنجی: درایه‌های پیشرو (ناصفر)؛ ستاره: هر عددی}}
\end{center}

\section{فرم پلکانی سطری‌-تحویل‌یافته (\lr{RREF})}

\begin{definition}[فرم پلکانی سطری‌-تحویل‌یافته]
ماتریسی که در فرم پلکانی است، در \term{فرم پلکانی سطری‌-تحویل‌یافته}{Reduced Row Echelon Form, RREF} است اگر علاوه بر آن:
\begin{enumerate}
  \item درایه‌ی پیشروی هر سطر ناصفر برابر ۱ باشد.
  \item هر ۱ِ پیشرو تنها درایه‌ی ناصفر ستون خودش باشد.
  \item ۱ِ پیشروی هر سطر (از سطر دوم به بعد) سمت راست ۱ِ پیشروی سطر بالایی‌اش باشد.
  \item سطرهای تمام‌صفر در پایین ماتریس باشند.
\end{enumerate}
به بیان کوتاه: \hl{RREF = فرم پلکانی + سطری‌-تحویل‌یافته.}
\end{definition}

مثلاً
\[
\begin{bmatrix}
1&0&0&29\\
0&1&0&16\\
0&0&1&3
\end{bmatrix}
\]
در فرم RREF است. اگر این ماتریس افزوده‌ی یک دستگاه باشد، جواب بی‌درنگ خوانده می‌شود: $x_1=29,\ x_2=16,\ x_3=3$. ستون‌های پیشرو همان بردارهای استاندارد $\ve_1,\ve_2,\ve_3$ هستند.

\begin{theorem}\label{thm:rref-exists}
هر ماتریس $m\times n$ با یک ماتریس در فرم RREF هم‌ارز سطری است.
\end{theorem}
\begin{proof}
طبق قضیه‌ی \ref{thm:row-reduced-exists}، $A$ با یک ماتریس سطری‌-تحویل‌یافته‌ی $R$ هم‌ارز سطری است. حالا فقط با عمل جابه‌جایی سطرها:
\begin{enumerate}
  \item سطرهای تمام‌صفر را به پایین ماتریس می‌بریم؛
  \item سطرهای ناصفر را طوری مرتب می‌کنیم که شماره‌ی ستون درایه‌ی پیشروی آن‌ها از بالا به پایین صعودی شود.
\end{enumerate}
جابه‌جایی سطرها، مقدار درایه‌های پیشرو (همه ۱) و صفر بودن بقیه‌ی درایه‌های ستون‌های پیشرو را خراب نمی‌کند؛ پس ماتریس هنوز سطری‌-تحویل‌یافته است و حالا شرط‌های پلکانی را هم دارد. در نتیجه در فرم RREF است. چون به $R$ فقط عمل سطری اضافه کردیم، این ماتریس با $A$ هم‌ارز سطری است.
\end{proof}

\begin{example}[RREF هست یا نه؟]
کدام ماتریس‌ها در فرم RREF هستند؟
\[
A=0_{m\times n},\quad
B=\begin{bmatrix}1&0&0\\0&2&-1\\0&0&0\end{bmatrix},\quad
C=\begin{bmatrix}0&1\\1&0\end{bmatrix},\quad
D=\begin{bmatrix}0&1&-3&0&0.5\\0&0&0&1&2\\0&0&0&0&0\end{bmatrix}.
\]
\end{example}
\begin{solution}
\begin{itemize}
  \item $A=0_{m\times n}$: \textcolor{thmcolor}{\textbf{بله.}} هیچ سطر ناصفری ندارد، پس همه‌ی شرط‌ها به‌طور بدیهی برقرارند.
  \item $B$: \textcolor{warncolor}{\textbf{خیر.}} درایه‌ی پیشروی سطر دوم ۲ است. (اصلاح: $R_2\leftarrow\frac12R_2$.)
  \item $C$: \textcolor{warncolor}{\textbf{خیر.}} درایه‌ی پیشروی سطر دوم (ستون ۱) سمت چپِ درایه‌ی پیشروی سطر اول (ستون ۲) است و این شرط پلکانی را نقض می‌کند. (اصلاح: $R_1\leftrightarrow R_2$.)
  \item $D$: \textcolor{thmcolor}{\textbf{بله.}} ۱ِ پیشروی سطر ۲ (ستون ۴) سمت راست ۱ِ پیشروی سطر ۱ (ستون ۲) است، سطر صفر پایین است و ستون‌های ۲ و ۴ غیر از ۱ِ پیشرو درایه‌ی ناصفری ندارند. (درایه‌های ستون‌های ۳ و ۵ می‌توانند هر عددی باشند.)
\end{itemize}
\end{solution}

\subsection{یکتایی RREF}
ممکن است یک ماتریس را با دنباله‌های متفاوتی از عملیات سطری کاهش دهیم. آیا همیشه به \emph{یک} RREF می‌رسیم؟ پاسخ مثبت است و همین باعث می‌شود بتوانیم از «\emph{همان} RREF ماتریس $A$» صحبت کنیم و آن را با $\RREF(A)$ نشان دهیم. اول یک لم کلیدی:

\begin{lemma}[عملیات سطری روابط بین ستون‌ها را حفظ می‌کنند]\label{lem:col-relations}
اگر $A$ و $B$ هم‌ارز سطری باشند، آن‌گاه برای هر بردار $\vx$:
\[
A\vx=\zero\iff B\vx=\zero .
\]
یعنی اگر بین ستون‌های $A$ رابطه‌ی خطی $x_1\va_1+\dots+x_n\va_n=\zero$ برقرار باشد، دقیقاً همان رابطه بین ستون‌های $B$ هم برقرار است و برعکس.
\end{lemma}
\begin{proof}
طبق قضیه‌ی \ref{thm:row-equiv-PA}، $B=PA$ که $P$ حاصل‌ضرب ماتریس‌های مقدماتی و در نتیجه وارون‌پذیر است (نتیجه‌ی \ref{cor:elem-invertible}). اگر $A\vx=\zero$ آن‌گاه $B\vx=P(A\vx)=P\zero=\zero$. برعکس اگر $B\vx=\zero$ آن‌گاه $A\vx=P^{-1}B\vx=\zero$.
\end{proof}

\begin{theorem}[یکتایی RREF]\label{thm:rref-unique}
هر ماتریس دقیقاً با \emph{یک} ماتریس در فرم RREF هم‌ارز سطری است.
\end{theorem}
\begin{proof}
فرض کنید $R$ در فرم RREF باشد و ستون‌های پیشروی آن $p_1<p_2<\dots<p_r$ باشند. چون RREF است، ستون $p_k$ از $R$ دقیقاً بردار استاندارد $\ve_k$ است. همچنین اگر ستون $j$ قبل از $p_{k+1}$ باشد، درایه‌های آن در سطرهای $k+1$ به بعد صفرند (چون درایه‌ی پیشروی آن سطرها دیرتر می‌آید). دو نتیجه می‌گیریم:
\begin{itemize}
  \item \textbf{(الف)} ستون $p_{k+1}=\ve_{k+1}$ ترکیب خطی ستون‌های قبل از خودش \emph{نیست}، چون همه‌ی آن ستون‌ها مؤلفه‌ی $(k+1)$-امشان صفر است. اما یک ستون غیرپیشروی $j$ (با $p_k<j<p_{k+1}$) ترکیب خطی ستون‌های قبلی هست:
  \[
  \text{\fa{ستون}}\ j=R_{1j}\ve_1+\dots+R_{kj}\ve_k=R_{1j}(\text{\fa{ستون}}\ p_1)+\dots+R_{kj}(\text{\fa{ستون}}\ p_k).
  \]
  \item \textbf{(ب)} پس «ستون $j$ پیشرو است» یعنی «ستون $j$ ترکیب خطی ستون‌های قبلی نیست»، و درایه‌های یک ستون غیرپیشرو همان ضرایبی هستند که آن ستون را بر حسب ستون‌های پیشرو می‌نویسند (این ضرایب یکتا هستند چون ستون‌های پیشرو بردارهای استاندارد متمایزند).
\end{itemize}
حالا فرض کنید $R$ و $S$ هر دو RREF و هم‌ارز سطری با $A$ باشند. پس $R$ و $S$ با هم هم‌ارز سطری‌اند و طبق لم \ref{lem:col-relations}، هر رابطه‌ی خطی بین ستون‌های $R$ دقیقاً بین ستون‌های $S$ هم برقرار است. «ستون $j$ ترکیبی از ستون‌های قبلی است» یک رابطه‌ی خطی است؛ پس طبق (الف) ستون‌های پیشروی $R$ و $S$ یکی هستند و در هر دو، ستون پیشروی $k$-ام برابر $\ve_k$ است. برای ستون‌های غیرپیشرو، طبق (ب) درایه‌ها همان ضرایب رابطه‌ی خطی‌اند که در هر دو ماتریس یکسان است. پس همه‌ی ستون‌های $R$ و $S$ برابرند: $R=S$.
\end{proof}

\begin{note}
فرم پلکانی (غیرتحویل‌یافته) یکتا نیست: مثلاً می‌توانیم هر سطرش را در ۲ ضرب کنیم و باز پلکانی بماند. اما \textbf{جای} درایه‌های پیشرو در همه‌ی فرم‌های پلکانی یک ماتریس یکسان است، چون هر فرم پلکانی با ادامه‌ی کاهش به همان RREF یکتا می‌رسد و این کار جای درایه‌های پیشرو را تغییر نمی‌دهد.
\end{note}

\subsection{مکان محوری و الگوریتم کاهش سطری}
\begin{definition}[مکان محوری و ستون محوری]
\term{مکان محوری}{Pivot Position} در ماتریس $A$ جایی است که در $\RREF(A)$ یک ۱ِ پیشرو قرار دارد. ستونی که مکان محوری دارد \term{ستون محوری}{Pivot Column} نامیده می‌شود و به درایه‌ی ناصفری که در آن مکان (در یک فرم پلکانی) می‌نشیند \term{محور}{Pivot} می‌گوییم.
\end{definition}

\begin{note}
در فرم RREF: \hl{تعداد سطرهای ناصفر = تعداد ستون‌های محوری}؛ چون هر سطر ناصفر دقیقاً یک ۱ِ پیشرو دارد و هر ۱ِ پیشرو در یک ستون متفاوت است.
\end{note}

\paragraph{الگوریتم کاهش سطری (حذف گاوسی-جردن).}
\begin{enumerate}
  \item \textbf{مرحله‌ی رفت} \enterm{Forward Phase}: چپ‌ترین ستون ناصفر را پیدا کن؛ این یک ستون محوری است. در صورت لزوم با جابه‌جایی، یک درایه‌ی ناصفر را به بالای آن ستون بیاور (محور). با عمل‌های جایگزینی، درایه‌های \emph{زیر} محور را صفر کن. سطر محوری را کنار بگذار و همین کار را روی زیرماتریس باقی‌مانده تکرار کن. نتیجه: \textbf{فرم پلکانی}.
  \item \textbf{مرحله‌ی برگشت} \enterm{Backward Phase}: از راست‌ترین محور شروع کن و با جایگزینی، درایه‌های \emph{بالای} هر محور را صفر کن و با مقیاس‌دهی هر محور را ۱ کن. نتیجه: \textbf{RREF}.
\end{enumerate}

\begin{example}[کاهش سطری کامل]
ماتریس زیر را به فرم RREF برسانید:
\[
A=\begin{bmatrix}0&3&-6&6&4&-5\\3&-7&8&-5&8&9\\3&-9&12&-9&6&15\end{bmatrix}.
\]
\end{example}
\begin{solution}
\textbf{مرحله‌ی رفت.} ستون اول، چپ‌ترین ستون ناصفر است اما بالایش صفر است؛ پس $R_1\leftrightarrow R_3$:
\[
\begin{bmatrix}\boxed{3}&-9&12&-9&6&15\\3&-7&8&-5&8&9\\0&3&-6&6&4&-5\end{bmatrix}
\xrightarrow{R_2\leftarrow R_2-R_1}
\begin{bmatrix}3&-9&12&-9&6&15\\0&\boxed{2}&-4&4&2&-6\\0&3&-6&6&4&-5\end{bmatrix}
\]
\[
\xrightarrow{R_3\leftarrow R_3-\frac32R_2}
\begin{bmatrix}3&-9&12&-9&6&15\\0&2&-4&4&2&-6\\0&0&0&0&\boxed{1}&4\end{bmatrix}.
\]
این فرم پلکانی است؛ ستون‌های محوری ۱، ۲ و ۵ هستند.

\textbf{مرحله‌ی برگشت.} اول درایه‌های بالای محور ستون ۵ را صفر می‌کنیم، بعد محورها را ۱ می‌کنیم:
\[
\xrightarrow[R_1\leftarrow R_1-6R_3]{R_2\leftarrow R_2-2R_3}
\begin{bmatrix}3&-9&12&-9&0&-9\\0&2&-4&4&0&-14\\0&0&0&0&1&4\end{bmatrix}
\xrightarrow{R_2\leftarrow\frac12R_2}
\begin{bmatrix}3&-9&12&-9&0&-9\\0&1&-2&2&0&-7\\0&0&0&0&1&4\end{bmatrix}
\]
\[
\xrightarrow{R_1\leftarrow R_1+9R_2}
\begin{bmatrix}3&0&-6&9&0&-72\\0&1&-2&2&0&-7\\0&0&0&0&1&4\end{bmatrix}
\xrightarrow{R_1\leftarrow\frac13R_1}
\begin{bmatrix}1&0&-2&3&0&-24\\0&1&-2&2&0&-7\\0&0&0&0&1&4\end{bmatrix}.
\]
این $\RREF(A)$ است.
\end{solution}

\section{جواب‌های یک دستگاه خطی}

\subsection{متغیرهای پایه‌ای و آزاد}
فرض کنید ماتریس افزوده‌ی یک دستگاه به این فرم RREF رسیده باشد:
\[
\left[\begin{array}{ccc|c}1&0&-5&1\\0&1&1&4\\0&0&0&0\end{array}\right]
\quad\Longleftrightarrow\quad
\begin{cases}x-5z=1\\ y+z=4\\ 0=0\end{cases}
\quad\Longrightarrow\quad
\begin{cases}x=1+5z\\ y=4-z\\ z\ \text{\fa{آزاد}}\end{cases}
\]
\begin{itemize}
  \item متغیرهایی که به ستون‌های محوری مربوط‌اند ($x$ و $y$) \term{متغیرهای پایه‌ای}{Basic Variables} هستند.
  \item بقیه ($z$) \term{متغیرهای آزاد}{Free Variables} هستند: هر مقداری می‌توانند بگیرند و متغیرهای پایه‌ای بر حسب آن‌ها تعیین می‌شوند.
\end{itemize}
این دستگاه سازگار است، چون توانستیم مجموعه‌ی جواب را صریحاً بنویسیم: به ازای هر انتخاب $z$ یک جواب داریم؛ پس بی‌شمار جواب.

\subsection{قضیه‌ی وجود و یکتایی}
\begin{theorem}[شرط سازگاری]\label{thm:consistency}
یک دستگاه خطی سازگار است اگر و تنها اگر راست‌ترین ستون ماتریس افزوده، ستون محوری \emph{نباشد}؛ یعنی اگر و تنها اگر در فرم پلکانی ماتریس افزوده، هیچ سطری به شکل
\[
\begin{bmatrix}0&0&\cdots&0&\big|&b\end{bmatrix}\qquad (b\neq0)
\]
وجود نداشته باشد. (مثلاً $\begin{bmatrix}0&\cdots&0&|&-2\end{bmatrix}$ یا $\begin{bmatrix}0&\cdots&0&|&20\end{bmatrix}$.)
\end{theorem}
\begin{proof}
طبق قضیه‌ی \ref{thm:row-equiv-solutions}، دستگاه اصلی و دستگاهی که ماتریس افزوده‌اش فرم پلکانی (یا RREF) آن است، مجموعه‌ی جواب یکسانی دارند؛ پس کافی است دستگاه کاهش‌یافته را بررسی کنیم.

\textbf{(رفت: اگر سازگار است، آخرین ستون محوری نیست.)} با برهان خلف: فرض کنید آخرین ستون محوری باشد. آن‌گاه محور آن ستون در سطری است که همه‌ی درایه‌های قبلی‌اش صفرند؛ یعنی سطری به شکل $\begin{bmatrix}0&\cdots&0&|&b\end{bmatrix}$ با $b\neq0$ داریم. این سطر معادله‌ی
\[
0x_1+0x_2+\dots+0x_n=b
\]
است که برای \emph{هیچ} انتخابی از $x_i$ها برقرار نیست (سمت چپ همیشه صفر است). پس دستگاه جواب ندارد؛ با فرض سازگاری تناقض دارد.

\textbf{(برگشت: اگر آخرین ستون محوری نیست، سازگار است.)} ماتریس افزوده را تا RREF کاهش دهید. چون آخرین ستون محوری نیست، هر سطر ناصفر یک ۱ِ پیشرو در یکی از ستون‌های \emph{ضرایب} دارد؛ یعنی هر معادله‌ی ناصفر به شکل
\[
x_{p}+(\text{\fa{ترکیبی از متغیرهای آزاد}})=d
\]
است و هر متغیر پایه‌ای فقط در یک معادله ظاهر می‌شود. سطرهای صفر هم معادله‌ی $0=0$ هستند که همیشه برقرارند. حالا به همه‌ی متغیرهای آزاد مقدار دلخواه (مثلاً صفر) بدهید و هر متغیر پایه‌ای را از معادله‌ی خودش حساب کنید. این یک جواب است؛ پس دستگاه سازگار است.
\end{proof}

\begin{theorem}[یکتایی]\label{thm:uniqueness}
اگر یک دستگاه خطی سازگار باشد، مجموعه‌ی جواب آن:
\begin{itemize}
  \item \textbf{دقیقاً یک جواب} دارد، اگر هیچ متغیر آزادی نباشد؛
  \item \textbf{بی‌شمار جواب} دارد، اگر دست‌کم یک متغیر آزاد باشد.
\end{itemize}
\end{theorem}
\begin{proof}
طبق اثبات قضیه‌ی قبل، هر جواب با انتخاب مقدار متغیرهای آزاد ساخته می‌شود و متغیرهای پایه‌ای به‌طور \emph{یکتا} از روی آن‌ها تعیین می‌شوند. اگر متغیر آزادی نباشد، هر متغیر یک مقدار مشخص دارد و فقط یک جواب هست. اگر متغیر آزادی مثل $x_j$ باشد، به ازای هر عدد $t$ با گذاشتن $x_j=t$ (و صفر برای بقیه‌ی آزادها) جوابی به دست می‌آید که مؤلفه‌ی $j$-امش $t$ است؛ پس جواب‌های متفاوت بی‌شماری داریم.
\end{proof}

\subsection{الگوریتم یافتن همه‌ی جواب‌ها}
\begin{enumerate}
  \item ماتریس افزوده‌ی دستگاه را بنویسید.
  \item با الگوریتم کاهش سطری، آن را به یک فرم پلکانی برسانید و سازگاری را بررسی کنید (قضیه‌ی \ref{thm:consistency}). اگر جوابی نیست، تمام.
  \item کاهش را تا RREF ادامه دهید.
  \item دستگاه معادلات متناظر با ماتریس مرحله‌ی ۳ را بنویسید.
  \item هر معادله‌ی ناصفر را طوری بازنویسی کنید که تنها متغیر پایه‌ای آن بر حسب متغیرهای آزاد بیان شود.
\end{enumerate}

\begin{example}
در مثال کاهش سطری کامل، فرض کنید $A$ ماتریس افزوده‌ی یک دستگاه با ۵ مجهول باشد. همه‌ی جواب‌ها را بنویسید.
\end{example}
\begin{solution}
از RREF به دست آمده:
\[
\begin{cases}x_1-2x_3+3x_4=-24\\ x_2-2x_3+2x_4=-7\\ x_5=4\end{cases}
\quad\Longrightarrow\quad
\begin{cases}x_1=-24+2x_3-3x_4\\ x_2=-7+2x_3-2x_4\\ x_3,\ x_4\ \text{\fa{آزاد}}\\ x_5=4\end{cases}
\]
آخرین ستون محوری نیست، پس دستگاه سازگار است و چون دو متغیر آزاد دارد، بی‌شمار جواب دارد.
\end{solution}

\begin{example}[سازگاری برای همه‌ی $\vb$ها]
فرض کنید
$A=\begin{bmatrix}1&3&4\\-4&2&-6\\-3&-2&-7\end{bmatrix}$ و $\vb=\begin{bmatrix}b_1\\b_2\\b_3\end{bmatrix}$.
آیا معادله‌ی $A\vx=\vb$ برای \emph{همه‌ی} مقادیر ممکن $b_1,b_2,b_3$ سازگار است؟
\end{example}
\begin{solution}
ماتریس افزوده را به فرم پلکانی می‌بریم:
\[
\left[\begin{array}{ccc|c}1&3&4&b_1\\-4&2&-6&b_2\\-3&-2&-7&b_3\end{array}\right]
\sim
\left[\begin{array}{ccc|c}1&3&4&b_1\\0&14&10&b_2+4b_1\\0&7&5&b_3+3b_1\end{array}\right]
\sim
\left[\begin{array}{ccc|c}1&3&4&b_1\\0&14&10&b_2+4b_1\\0&0&0&b_3+3b_1-\frac12(b_2+4b_1)\end{array}\right].
\]
درایه‌ی سوم ستون آخر برابر $b_1-\frac12b_2+b_3$ است. طبق قضیه‌ی \ref{thm:consistency}، دستگاه فقط وقتی سازگار است که این عدد صفر باشد. پس $A\vx=\vb$ \textbf{برای همه‌ی $\vb$ها سازگار نیست}؛ مثلاً برای $\vb=(1,0,0)$ داریم $b_1-\frac12b_2+b_3=1\neq0$. در واقع فقط بردارهایی که روی صفحه‌ی $b_1-\frac12b_2+b_3=0$ هستند جواب دارند.
\end{solution}

مشکل این مثال این بود که $A$ در سطر سوم محور نداشت. قضیه‌ی زیر این را دقیق می‌کند:

\begin{theorem}[سازگاری برای هر $\vb$]\label{thm:pivot-every-row}
برای ماتریس $m\times n$ مثل $A$ این دو گزاره هم‌ارزند:
\begin{enumerate}
  \item معادله‌ی $A\vx=\vb$ برای \emph{هر} $\vb\in\R^m$ دست‌کم یک جواب دارد.
  \item ماتریس ضرایب $A$ در \emph{هر سطر} یک مکان محوری دارد.
\end{enumerate}
\end{theorem}
\begin{proof}
فرض کنید $U$ یک فرم پلکانی $A$ باشد و $U=PA$ با $P$ وارون‌پذیر.

\textbf{(از ۲ به ۱)} اگر هر سطر $U$ محور داشته باشد، در فرم پلکانی $\aug{A}{\vb}$ (که همان $\aug{U}{P\vb}$ است) هیچ سطری به شکل $\begin{bmatrix}0&\cdots&0&|&b\end{bmatrix}$ نمی‌تواند باشد، چون هر سطر یک محور در بخش ضرایب دارد. پس آخرین ستون محوری نیست و طبق قضیه‌ی \ref{thm:consistency} دستگاه سازگار است.

\textbf{(از ۱ به ۲)} با برهان خلف فرض کنید سطر آخر $U$ محور نداشته باشد؛ پس سطر آخر $U$ تمام‌صفر است. بردار $\vb=P^{-1}\ve_m$ را در نظر بگیرید. عملیات سطری که $A$ را به $U$ می‌برد، $\aug{A}{\vb}$ را به $\aug{PA}{P\vb}=\aug{U}{\ve_m}$ می‌برد که سطر آخرش $\begin{bmatrix}0&\cdots&0&|&1\end{bmatrix}$ است. پس برای این $\vb$ دستگاه ناسازگار است؛ تناقض.
\end{proof}

\begin{example}[درست یا نادرست؟]
«معادله‌ی $A\vx=\vb$ سازگار است اگر ماتریس افزوده‌ی $\aug{A}{\vb}$ در فرم RREF در هر سطر یک ستون محوری داشته باشد.»
\end{example}
\begin{solution}
\textbf{نادرست.} محور یک سطر ممکن است در \emph{آخرین} ستون (ستون $\vb$) باشد. مثال نقض:
\[
\aug{A}{\vb}=\left[\begin{array}{cc|c}1&0&0\\0&0&1\end{array}\right].
\]
هر دو سطر محور دارند، اما سطر دوم یعنی $0=1$؛ دستگاه ناسازگار است. نسخه‌ی درست جمله این است: «اگر \textbf{ماتریس ضرایب} $A$ در هر سطر محور داشته باشد، $A\vx=\vb$ سازگار است» (قضیه‌ی \ref{thm:pivot-every-row}).
\end{solution}

\subsection{دستگاه‌های همگن}
\begin{definition}[دستگاه همگن]
دستگاهی که بتوان آن را به شکل $A\vx=\zero$ نوشت \textbf{همگن} است.
\begin{itemize}
  \item \term{جواب بدیهی}{Trivial Solution}: $A\vx=\zero$ همیشه دست‌کم یک جواب دارد: $\vx=\zero$.
  \item \term{جواب نابدیهی}{Nontrivial Solution}: هر جواب ناصفر $A\vx=\zero$.
\end{itemize}
\end{definition}
پس دستگاه همگن \emph{همیشه} سازگار است و سؤال اصلی فقط یکتایی است.

\begin{theorem}\label{thm:homog-free}
معادله‌ی همگن $A\vx=\zero$ جواب نابدیهی دارد اگر و تنها اگر دست‌کم یک متغیر آزاد داشته باشد.
\end{theorem}
\begin{proof}
فرض کنید $A$ ماتریس $m\times n$ باشد و آن را به RREF برسانیم. دو حالت داریم:
\begin{enumerate}
  \item \textbf{همه‌ی متغیرها پایه‌ای‌اند:} هر متغیر، متغیر پیشروی یک سطر است و معادله‌ی آن سطر می‌گوید $x_j=0$ (سمت راست صفر است و متغیر آزادی نیست). پس تنها جواب $\vx=\zero$ است.
  \item \textbf{دست‌کم یک متغیر آزاد هست:} فرض کنید $x_j$ آزاد باشد. بگذارید $x_j=t\neq0$ (و بقیه‌ی آزادها صفر) و متغیرهای پایه‌ای را حساب کنید. جوابی به دست می‌آید که مؤلفه‌ی $j$-امش $t\neq0$ است؛ پس یک جواب نابدیهی داریم.
\end{enumerate}
\end{proof}

\begin{example}
\[
\left[\begin{array}{ccc|c}1&3&4&0\\2&-1&2&0\\1&0&1&0\end{array}\right]
\ \xrightarrow{\RREF}\
\left[\begin{array}{ccc|c}1&0&0&0\\0&1&0&0\\0&0&1&0\end{array}\right]
\]
هیچ متغیر آزادی نیست؛ پس فقط جواب بدیهی دارد. اما دستگاه
$\left[\begin{array}{cccc|c}1&0&-8&-7&0\\0&1&4&3&0\\0&0&0&0&0\end{array}\right]$
دو متغیر آزاد ($x_3,x_4$) دارد و بی‌شمار جواب نابدیهی.
\end{example}

\begin{corollary}[مجهول‌های بیشتر از معادله‌ها]\label{cor:more-unknowns}
اگر یک دستگاه همگن تعداد مجهول‌هایش از تعداد معادله‌هایش بیشتر باشد ($n>m$)، جواب نابدیهی دارد.
\end{corollary}
\begin{proof}
هر محور در یک سطر متفاوت است، پس حداکثر $m$ محور (و $m$ متغیر پایه‌ای) داریم. چون $n>m$، دست‌کم $n-m\ge1$ متغیر آزاد وجود دارد و طبق قضیه‌ی \ref{thm:homog-free} جواب نابدیهی هست.
\end{proof}

\begin{note}
این نتیجه‌ی به‌ظاهر ساده، در فصل ۶ کلید اثبات مهم‌ترین قضیه‌ی بُعد است.
\end{note}

\subsection{وجود جواب و ترکیب خطی ستون‌ها}
\begin{theorem}\label{thm:solvable-combination}
معادله‌ی $A\vx=\vb$ جواب دارد اگر و تنها اگر $\vb$ ترکیبی خطی از ستون‌های $A$ باشد.
\end{theorem}
\begin{proof}
فرض کنید $A\in\R^{m\times n}$ ستون‌های $\va_1,\dots,\va_n$ را داشته باشد. طبق نگاه ستونی ضرب ماتریس در بردار، برای هر $\vx$:
\[
A\vx=x_1\va_1+x_2\va_2+\dots+x_n\va_n .
\]
\textbf{(رفت)} اگر $A\vx=\vb$ جوابی مثل $\vx$ داشته باشد، $\vb=x_1\va_1+\dots+x_n\va_n$؛ یعنی $\vb$ ترکیب خطی ستون‌هاست.

\textbf{(برگشت)} اگر $\vb=c_1\va_1+\dots+c_n\va_n$ برای اعدادی مثل $c_i$، بگذارید $\vx=(c_1,\dots,c_n)\T$. آن‌گاه $A\vx=c_1\va_1+\dots+c_n\va_n=\vb$؛ پس $A\vx=\vb$ جواب دارد.
\end{proof}

\subsection{ساختار جواب دستگاه ناهمگن}
\begin{theorem}[جواب عمومی = جواب خصوصی + جواب همگن]\label{thm:general-solution}
فرض کنید $A\vx=\vb$ سازگار باشد و $\vp$ یک جواب آن باشد. آن‌گاه مجموعه‌ی همه‌ی جواب‌های $A\vx=\vb$ دقیقاً برابر است با
\[
\{\vp+\vv_h\ :\ A\vv_h=\zero\}.
\]
\end{theorem}
\begin{proof}
\textbf{(هر $\vp+\vv_h$ جواب است)} $A(\vp+\vv_h)=A\vp+A\vv_h=\vb+\zero=\vb$.

\textbf{(هر جواب این شکلی است)} اگر $A\vx=\vb$، بگذارید $\vv_h=\vx-\vp$. آن‌گاه $A\vv_h=A\vx-A\vp=\vb-\vb=\zero$ و $\vx=\vp+\vv_h$.
\end{proof}

\begin{intuition}
مجموعه‌ی جواب $A\vx=\vb$ همان مجموعه‌ی جواب $A\vx=\zero$ است که به اندازه‌ی $\vp$ \textbf{جابه‌جا} شده است. اگر جواب‌های همگن یک خط گذرنده از مبدأ باشند، جواب‌های ناهمگن خطی موازی با آن هستند.
\begin{center}
\begin{tikzpicture}[scale=0.8,>=Latex]
  \draw[->] (-2.5,0)--(4,0); \draw[->] (0,-1.2)--(0,3.4);
  \draw[very thick, primary] (-2,-1)--(3.6,1.8) node[right] {\small $A\vx=\zero$};
  \draw[very thick, excolor] (-2.4,0.8)--(3,3.5) node[right] {\small $A\vx=\vb$};
  \draw[->, thick, thmcolor] (0,0)--(-0.4,1.6) node[left] {\small $\vp$};
  \draw[->, thick, primary!70] (0,0)--(2,1) node[below right] {\small $\vv_h$};
  \draw[->, thick, warncolor] (0,0)--(1.6,2.6) node[above] {\small $\vp+\vv_h$};
  \draw[dashed, gray] (2,1)--(1.6,2.6); \draw[dashed, gray] (-0.4,1.6)--(1.6,2.6);
\end{tikzpicture}
\end{center}
\end{intuition}

\section{تجزیه‌ی $LU$}

\subsection{چرا تجزیه؟}
همان‌طور که تجزیه‌ی یک عدد طبیعی به حاصل‌ضرب اعداد دیگر (مثلاً $30=2\times3\times5$) طبیعی و روشنگر است، نوشتن یک ماتریس به صورت حاصل‌ضرب چند ماتریس دیگر هم اغلب مفید است. به چنین نمایشی \term{تجزیه‌ی ماتریسی}{Matrix Factorization} می‌گوییم. همه‌ی تجزیه‌ها به یک اندازه مفید نیستند؛ در این بخش تجزیه‌ای را می‌بینیم که از دل حذف گاوسی بیرون می‌آید و برای پیاده‌سازی در کامپیوتر بسیار مناسب است. مثلاً
\[
\begin{bmatrix}3&-1\\9&-5\end{bmatrix}=\begin{bmatrix}1&0\\3&1\end{bmatrix}\begin{bmatrix}3&-1\\0&-2\end{bmatrix}.
\]

\subsection{حذف گاوسی به زبان ماتریس‌های مقدماتی}
دستگاه $A\vx=\vb$ با $A$ مربعی $n\times n$ را در نظر بگیرید. می‌خواهیم نشان دهیم حذف گاوسی، به‌طور ضمنی، $A$ را به حاصل‌ضرب دو ماتریس تجزیه می‌کند که با آن‌ها حل این دستگاه (و هر دستگاه دیگری با همین ماتریس ضرایب) آسان می‌شود.

\begin{example}
$A=\begin{bmatrix}2&1&3\\4&-1&3\\-2&5&5\end{bmatrix}$ را با حذف گاوسی به فرم پلکانی ببرید و ماتریس‌های مقدماتی را ثبت کنید.
\end{example}
\begin{solution}
\[
A=\begin{bmatrix}2&1&3\\4&-1&3\\-2&5&5\end{bmatrix}
\xrightarrow{R_2\leftarrow R_2-2R_1}
\begin{bmatrix}2&1&3\\0&-3&-3\\-2&5&5\end{bmatrix}
\]
\[
\xrightarrow{R_3\leftarrow R_3+R_1}
\begin{bmatrix}2&1&3\\0&-3&-3\\0&6&8\end{bmatrix}
\xrightarrow{R_3\leftarrow R_3+2R_2}
\begin{bmatrix}2&1&3\\0&-3&-3\\0&0&2\end{bmatrix}=U .
\]
ماتریس‌های مقدماتی این سه عمل (به ترتیب):
\[
E_1=\begin{bmatrix}1&0&0\\-2&1&0\\0&0&1\end{bmatrix},\quad
E_2=\begin{bmatrix}1&0&0\\0&1&0\\1&0&1\end{bmatrix},\quad
E_3=\begin{bmatrix}1&0&0\\0&1&0\\0&2&1\end{bmatrix}.
\]
پس $E_3E_2E_1A=U$. با ضرب دو طرف از چپ در $E_3^{-1}$، سپس $E_2^{-1}$ و سپس $E_1^{-1}$:
\[
A=E_1^{-1}E_2^{-1}E_3^{-1}U=
\begin{bmatrix}1&0&0\\2&1&0\\0&0&1\end{bmatrix}
\begin{bmatrix}1&0&0\\0&1&0\\-1&0&1\end{bmatrix}
\begin{bmatrix}1&0&0\\0&1&0\\0&-2&1\end{bmatrix}U
=\underbrace{\begin{bmatrix}1&0&0\\2&1&0\\-1&-2&1\end{bmatrix}}_{L}U .
\]
پس $A=LU$.
\end{solution}

\begin{definition}[تجزیه‌ی $LU$]
فرض کنید $A$ ماتریسی مربعی باشد. نمایش $A$ به صورت $A=LU$ که در آن $L$ \term{پایین‌مثلثی واحد}{Unit Lower Triangular} و $U$ بالامثلثی است، \term{تجزیه‌ی $LU$}{LU Factorization} ماتریس $A$ نامیده می‌شود. «پایین‌مثلثی واحد» یعنی بالای قطر صفر و روی قطر همه ۱:
\[
L=\begin{bmatrix}1&0&\cdots&0\\ *&1&\cdots&0\\ \vdots&\vdots&\ddots&\vdots\\ *&*&\cdots&1\end{bmatrix}.
\]
\end{definition}

\subsection{شرط وجود تجزیه‌ی $LU$}
\begin{lemma}\label{lem:unit-lower}
\begin{enumerate}
  \item حاصل‌ضرب دو ماتریس پایین‌مثلثی واحد (هم‌اندازه)، پایین‌مثلثی واحد است.
  \item ماتریس مقدماتی عمل $R_i\leftarrow R_i+cR_j$ با $j<i$ (اضافه کردن مضربی از یک سطر \emph{بالاتر} به یک سطر \emph{پایین‌تر}) پایین‌مثلثی واحد است و وارونش (عمل $R_i\leftarrow R_i-cR_j$) هم همین‌طور.
\end{enumerate}
\end{lemma}
\begin{proof}
\textbf{۱.} فرض کنید $L,M$ پایین‌مثلثی واحد باشند؛ یعنی $L_{ik}=0$ وقتی $k>i$ و $M_{kj}=0$ وقتی $k<j$، و درایه‌های قطری ۱ هستند. در
$(LM)_{ij}=\sum_k L_{ik}M_{kj}$
یک جمله فقط وقتی می‌تواند ناصفر باشد که $j\le k\le i$.
\begin{itemize}
  \item اگر $i<j$، هیچ $k$ای در این بازه نیست، پس $(LM)_{ij}=0$: بالای قطر صفر است.
  \item اگر $i=j$، فقط $k=i$ می‌ماند: $(LM)_{ii}=L_{ii}M_{ii}=1$.
\end{itemize}
\textbf{۲.} این ماتریس همان $I$ است با یک درایه‌ی اضافه‌ی $c$ در جای $(i,j)$؛ چون $j<i$ این درایه زیر قطر است. وارون آن هم $I$ با $-c$ در همان جاست.
\end{proof}

\begin{theorem}[وجود تجزیه‌ی $LU$]\label{thm:lu-exists}
اگر ماتریس مربعی $A$ را بتوان \textbf{بدون جابه‌جایی سطرها} به فرم پلکانی رساند، $A$ تجزیه‌ی $LU$ دارد.
\end{theorem}
\begin{proof}
در مرحله‌ی رفتِ حذف گاوسی، اگر به جابه‌جایی نیاز نباشد، فقط از عمل‌های جایگزینی $R_i\leftarrow R_i-\ell_{ij}R_j$ با $j<i$ استفاده می‌کنیم (برای رسیدن به فرم پلکانی به مقیاس‌دهی هم نیازی نیست). پس
\[
E_k\cdots E_2E_1A=U
\]
که $U$ پلکانی و در نتیجه بالامثلثی است و هر $E_t$ پایین‌مثلثی واحد است. هر $E_t$ وارون‌پذیر است، پس
\[
A=(E_k\cdots E_1)^{-1}U=E_1^{-1}E_2^{-1}\cdots E_k^{-1}U .
\]
طبق لم \ref{lem:unit-lower}، هر $E_t^{-1}$ پایین‌مثلثی واحد است و حاصل‌ضربشان هم پایین‌مثلثی واحد است. پس $L:=E_1^{-1}\cdots E_k^{-1}$ پایین‌مثلثی واحد است و $A=LU$.
\end{proof}

\begin{warning}[وقتی جابه‌جایی لازم است]
اگر در حین حذف، در مکان محوری صفر بیاید و مجبور به جابه‌جایی سطر شویم، ماتریس مقدماتی جابه‌جایی پایین‌مثلثی نیست و $L$ دیگر پایین‌مثلثی واحد نمی‌شود. مثلاً $A=\begin{bmatrix}0&1\\1&0\end{bmatrix}$ تجزیه‌ی $LU$ ندارد: اگر
$\begin{bmatrix}1&0\\ \ell&1\end{bmatrix}\begin{bmatrix}a&b\\0&d\end{bmatrix}=\begin{bmatrix}a&b\\ \ell a&\ell b+d\end{bmatrix}=\begin{bmatrix}0&1\\1&0\end{bmatrix}$
آن‌گاه $a=0$ و هم‌زمان $\ell a=1$، که ممکن نیست. در چنین حالتی از تجزیه‌ی $PA=LU$ استفاده می‌شود که $P$ ماتریس جای‌گشت (حاصل‌ضرب ماتریس‌های جابه‌جایی) است.
\end{warning}

\subsection{راه سریع ساختن $L$}
لازم نیست ماتریس‌های مقدماتی را بنویسیم و در هم ضرب کنیم. قضیه‌ی زیر می‌گوید $L$ را می‌شود مستقیم از \textbf{ضریب‌هایی} که در حذف استفاده کردیم پر کرد.

\begin{theorem}[ضرب‌کننده‌ها درایه‌های $L$ هستند]\label{thm:lu-multipliers}
فرض کنید حذف گاوسی بدون جابه‌جایی، ستون به ستون (از چپ به راست) انجام شود و برای صفر کردن درایه‌ی $(i,j)$ از عمل $R_i\leftarrow R_i-\ell_{ij}R_j$ استفاده شود. آن‌گاه
\[
L=\begin{bmatrix}1&&&\\ \ell_{21}&1&&\\ \ell_{31}&\ell_{32}&1&\\ \vdots&\vdots&&\ddots\end{bmatrix};
\]
یعنی درایه‌ی $(i,j)$ ماتریس $L$ دقیقاً ضرب‌کننده‌ی $\ell_{ij}$ است و روی قطر ۱ قرار دارد.
\end{theorem}
\begin{proof}
وارون عمل $R_i\leftarrow R_i-\ell_{ij}R_j$ ماتریس $I+\ell_{ij}\ve_i\ve_j\T$ است (همانی با $\ell_{ij}$ در جای $(i,j)$). داریم $L=E_1^{-1}E_2^{-1}\cdots E_k^{-1}$. این حاصل‌ضرب را از چپ به راست می‌سازیم: از $M=I$ شروع می‌کنیم و هر بار $M$ را از راست در $E_t^{-1}$ ضرب می‌کنیم. اما
\[
M\,(I+\ell_{ij}\ve_i\ve_j\T)=M+\ell_{ij}(M\ve_i)\ve_j\T ,
\]
یعنی این ضرب، $\ell_{ij}$ برابرِ \textbf{ستون $i$-ام} $M$ را به \textbf{ستون $j$-ام} $M$ اضافه می‌کند. ادعا می‌کنیم در این لحظه ستون $i$-ام $M$ هنوز همان $\ve_i$ است: ستون $i$ فقط با عمل‌هایی تغییر می‌کند که «ستون مقصدشان» $i$ است، یعنی عمل‌های مربوط به حذف زیر محورِ ستون $i$؛ و چون حذف ستون به ستون انجام می‌شود و $j<i$، این عمل‌ها \emph{بعد از} همه‌ی عمل‌های ستون $j$ می‌آیند. پس این ضرب دقیقاً $\ell_{ij}\ve_i$ را به ستون $j$ اضافه می‌کند؛ یعنی عدد $\ell_{ij}$ را در جای $(i,j)$ می‌نشاند و به هیچ درایه‌ی دیگری دست نمی‌زند. در پایان، $L$ همان $I$ است که ضرب‌کننده‌ها در جای خودشان نشسته‌اند.
\end{proof}

\begin{note}[دستور عملی]
برای هر عمل $R_i\leftarrow R_i-k\,R_j$ در حذف گاوسی، عدد $k$ (با همین علامت) در درایه‌ی $(i,j)$ ماتریس $L$ نوشته می‌شود و قطر $L$ تمام ۱ است. اگر عمل به شکل $R_i\leftarrow R_i+cR_j$ نوشته شده، در $L$ عدد $-c$ را بنویسید. در مثال قبل: $R_2-2R_1$ و $R_3-(-1)R_1$ و $R_3-(-2)R_2$؛ پس $\ell_{21}=2,\ \ell_{31}=-1,\ \ell_{32}=-2$، که همان $L$ به دست آمده است.
\end{note}

\begin{example}[ساختن تجزیه‌ی $LU$]
تجزیه‌ی $LU$ ماتریس $A=\begin{bmatrix}6&-2&0\\9&-1&1\\3&7&5\end{bmatrix}$ را بیابید.
\end{example}
\begin{solution}
$L$ را با علامت $\bullet$ برای درایه‌های ناشناخته شروع می‌کنیم و حذف را انجام می‌دهیم:
\[
A=\begin{bmatrix}6&-2&0\\9&-1&1\\3&7&5\end{bmatrix}
\xrightarrow[R_3\leftarrow R_3-\frac36R_1]{R_2\leftarrow R_2-\frac96R_1}
\begin{bmatrix}6&-2&0\\0&2&1\\0&8&5\end{bmatrix}
\xrightarrow{R_3\leftarrow R_3-\frac82R_2}
\begin{bmatrix}6&-2&0\\0&2&1\\0&0&1\end{bmatrix}=U,
\]
و هم‌زمان ضرب‌کننده‌ها را در $L$ می‌نشانیم:
\[
L:\ \begin{bmatrix}1&0&0\\ \bullet&1&0\\ \bullet&\bullet&1\end{bmatrix}\to
\begin{bmatrix}1&0&0\\ \frac32&1&0\\ \frac12&\bullet&1\end{bmatrix}\to
\begin{bmatrix}1&0&0\\ \frac32&1&0\\ \frac12&4&1\end{bmatrix}.
\]
پس
\[
A=LU=\begin{bmatrix}1&0&0\\ \frac32&1&0\\ \frac12&4&1\end{bmatrix}\begin{bmatrix}6&-2&0\\0&2&1\\0&0&1\end{bmatrix}.
\]
(با ضرب کردن $L$ در $U$ درستی را بررسی کنید.)
\end{solution}

\subsection{تجزیه‌ی $LU$ برای ماتریس‌های غیرمربعی}
ایده‌ی $LU$ را می‌توان به ماتریس‌های غیرمربعی $m\times n$ هم گسترش داد: $L$ مربعی $m\times m$ و پایین‌مثلثی واحد است و $U$ یک ماتریس $m\times n$ در \textbf{فرم پلکانی}:
\[
\underbrace{A}_{3\times4}=\underbrace{L}_{3\times3}\cdot\underbrace{U}_{3\times4}:\qquad
\begin{bmatrix}*&0&0\\ *&*&0\\ *&*&*\end{bmatrix}
\begin{bmatrix}*&*&*&*\\0&*&*&*\\0&0&*&*\end{bmatrix}
=\begin{bmatrix}*&*&*&*\\ *&*&*&*\\ *&*&*&*\end{bmatrix}.
\]

\begin{example}
تجزیه‌ی $LU$ ماتریس $A=\begin{bmatrix}2&4&2&1\\1&3&1&1\\3&7&4&2\end{bmatrix}$ را بیابید.
\end{example}
\begin{solution}
\[
\xrightarrow[R_3\leftarrow R_3-\frac32R_1]{R_2\leftarrow R_2-\frac12R_1}
\begin{bmatrix}2&4&2&1\\0&1&0&0.5\\0&1&1&0.5\end{bmatrix}
\xrightarrow{R_3\leftarrow R_3-1\cdot R_2}
\begin{bmatrix}2&4&2&1\\0&1&0&0.5\\0&0&1&0\end{bmatrix}=U .
\]
ضرب‌کننده‌ها: $\ell_{21}=\frac12$، $\ell_{31}=\frac32$، $\ell_{32}=1$. پس
\[
L=\begin{bmatrix}1&0&0\\ \frac12&1&0\\ \frac32&1&1\end{bmatrix},\qquad
U=\begin{bmatrix}2&4&2&1\\0&1&0&0.5\\0&0&1&0\end{bmatrix}.
\]
\end{solution}

\subsection{حل دستگاه با تجزیه‌ی $LU$}
\begin{keyformula}
\[
A\vx=\vb\iff L\underbrace{(U\vx)}_{\vy}=\vb\iff
\begin{cases}L\vy=\vb & \text{\fa{(جایگذاری پیشرو)}}\\ U\vx=\vy & \text{\fa{(جایگذاری پسرو)}}\end{cases}
\]
\end{keyformula}
\begin{enumerate}
  \item دستگاه $A\vx=\vb$ را به شکل $LU\vx=\vb$ بنویسید.
  \item بردار جدید $\vy$ (با اندازه‌ی $n\times1$) را با $U\vx=\vy$ تعریف کنید.
  \item با این تعریف، $LU\vx=\vb$ می‌شود $L\vy=\vb$؛ آن را برای $\vy$ حل کنید (چون $L$ پایین‌مثلثی است، از بالا به پایین: \term{جایگذاری پیشرو}{Forward Substitution}).
  \item $\vy$ را در $U\vx=\vy$ بگذارید و $\vx$ را حل کنید (از پایین به بالا: جایگذاری پسرو).
\end{enumerate}

\begin{example}
با استفاده از تجزیه‌ی $LU$ ماتریس $A=\begin{bmatrix}2&1&3\\4&-1&3\\-2&5&5\end{bmatrix}$، دستگاه $A\vx=\vb$ را با $\vb=\begin{bmatrix}1\\-4\\9\end{bmatrix}$ حل کنید.
\end{example}
\begin{solution}
قبلاً دیدیم
$A=\begin{bmatrix}1&0&0\\2&1&0\\-1&-2&1\end{bmatrix}\begin{bmatrix}2&1&3\\0&-3&-3\\0&0&2\end{bmatrix}=LU$.

\textbf{گام اول: $L\vy=\vb$.}
\[
\begin{cases}y_1=1\\ 2y_1+y_2=-4\\ -y_1-2y_2+y_3=9\end{cases}
\Longrightarrow\ y_1=1,\quad y_2=-4-2y_1=-6,\quad y_3=9+y_1+2y_2=-2 .
\]
\textbf{گام دوم: $U\vx=\vy$.}
\[
\begin{cases}2x_1+x_2+3x_3=1\\ -3x_2-3x_3=-6\\ 2x_3=-2\end{cases}
\Longrightarrow\ 
\begin{cases}x_3=-1\\ -3x_2=-6+3x_3=-9\ \Rightarrow\ x_2=3\\ 2x_1=1-x_2-3x_3=1\ \Rightarrow\ x_1=\tfrac12\end{cases}
\]
پس جواب $\vx=\begin{bmatrix}\frac12\\3\\-1\end{bmatrix}$ است.
\end{solution}

\begin{note}[چرا $LU$ مفید است؟]
خودِ تجزیه‌ی یک ماتریس $n\times n$ حدود $\frac23n^3$ عمل حسابی هزینه دارد؛ اما بعد از آن، حل $L\vy=\vb$ و $U\vx=\vy$ روی هم فقط حدود $2n^2$ عمل می‌خواهد. پس اگر بخواهیم دستگاه‌هایی با \emph{همان} $A$ و $\vb$های متفاوت را حل کنیم (که در مهندسی بسیار پیش می‌آید)، کافی است یک بار تجزیه کنیم.
\end{note}

\section{تعبیر هندسی}

\subsection{خط در $\R^2$}
خط $\ell$ با معادله‌ی $2x+y=0$ را در نظر بگیرید. اگر $\vn=\begin{bmatrix}2\\1\end{bmatrix}$ و $\vx=\begin{bmatrix}x\\y\end{bmatrix}$، معادله می‌شود $\vn\cdot\vx=0$؛ یعنی خط، مجموعه‌ی بردارهایی است که بر $\vn$ عمودند. به $\vn$ \term{بردار نرمال}{Normal Vector} خط می‌گوییم. همچنین با بردار جهت $\vd=\begin{bmatrix}1\\-2\end{bmatrix}$ (که بر $\vn$ عمود است) می‌توان نوشت $\vx=t\vd$.

\begin{center}
\begin{tikzpicture}[scale=0.7,>=Latex]
  \draw[->] (-2.5,0)--(3,0) node[below] {$x$}; \draw[->] (0,-3)--(0,3) node[left] {$y$};
  \draw[very thick, primary] (-1.4,2.8)--(1.4,-2.8) node[right] {$\ell$};
  \draw[->, thick, excolor] (0,0)--(2,1) node[right] {$\vn$};
  \draw[->, thick, thmcolor] (0,0)--(0.8,-1.6) node[right] {$\vd$};
  \draw (0.22,-0.44)--(0.66,-0.22)--(0.44,0.22);
  \begin{scope}[xshift=7.5cm]
    \draw[->] (-2,0)--(3.5,0); \draw[->] (0,-1)--(0,3);
    \draw[very thick, primary] (-1.5,2.5)--(3,-0.5) node[right] {$\ell$};
    \fill (0.75,1) circle (2pt) node[below left] {$\vp$};
    \draw[->, thick, excolor] (0.75,1)--(1.55,2.2) node[above] {$\vn$};
    \fill (2.25,0) circle (2pt) node[below] {$\vx$};
    \draw[->, thick, warncolor] (0.75,1)--(2.2,0.04);
    \node[warncolor] at (1.9,0.9) {\small $\vx-\vp$};
    \draw[->, gray] (0,0)--(0.72,0.96);
  \end{scope}
\end{tikzpicture}
\end{center}

\begin{definition}[شکل‌های معادله‌ی خط در $\R^2$]
اگر $\vp$ نقطه‌ای مشخص روی خط $\ell$ و $\vn\neq\zero$ بردار نرمال $\ell$ باشد:
\begin{itemize}
  \item \textbf{شکل نرمال:} $\vn\cdot(\vx-\vp)=0$ یا $\vn\cdot\vx=\vn\cdot\vp$.
  \item \textbf{شکل عمومی:} $ax+by=c$ که $\vn=\begin{bmatrix}a\\b\end{bmatrix}$ بردار نرمال است.
  \item \textbf{شکل برداری:} $\vx=\vp+t\vd$ که $\vd\neq\zero$ بردار جهت خط است و $t\in\R$.
  \item \textbf{شکل پارامتری:} $\begin{cases}x=p_1+td_1\\ y=p_2+td_2\end{cases}$.
\end{itemize}
\end{definition}

\subsection{خط و صفحه در $\R^3$}
در $\R^3$ با $\vn=\begin{bmatrix}a\\b\\c\end{bmatrix}$ و $\vx=\begin{bmatrix}x\\y\\z\end{bmatrix}$، معادله‌ی $\vn\cdot(\vx-\vp)=0$ دیگر یک خط نیست بلکه یک \textbf{صفحه} است: $ax+by+cz=d$ که $d=\vn\cdot\vp$. صفحه را می‌توان با یک نقطه و دو بردار جهت (غیرهم‌راستا) $\vu,\vv$ هم توصیف کرد: $\vx-\vp=s\vu+t\vv$.

\begin{center}
\begin{tikzpicture}[scale=0.9,>=Latex]
  \fill[primary!12] (0,0)--(4,0.6)--(5,2.4)--(1,1.8)--cycle;
  \draw[primary!50] (0,0)--(4,0.6)--(5,2.4)--(1,1.8)--cycle;
  \fill (1.6,0.9) circle (2pt) node[below left] {$P$};
  \draw[->, thick, excolor] (1.6,0.9)--(1.6,2.7) node[above] {$\vn$};
  \fill (3.6,1.4) circle (2pt) node[right] {$X$};
  \draw[->, thick, warncolor] (1.6,0.9)--(3.55,1.38) node[midway, below] {\small $\vx-\vp$};
  \node at (2.5,-0.6) {\small $\vn\cdot(\vx-\vp)=0$};
  \begin{scope}[xshift=7cm]
    \fill[primary!12] (0,0)--(4,0.6)--(5,2.4)--(1,1.8)--cycle;
    \draw[primary!50] (0,0)--(4,0.6)--(5,2.4)--(1,1.8)--cycle;
    \fill (1.2,0.8) circle (2pt) node[left] {$P$};
    \draw[->, thick, thmcolor] (1.2,0.8)--(2.4,0.95) node[below] {$\vu$};
    \draw[->, thick, thmcolor] (1.2,0.8)--(1.65,1.6) node[left] {$\vv$};
    \draw[->, thick, warncolor] (1.2,0.8)--(3.9,1.9);
    \fill (3.9,1.9) circle (2pt) node[right] {$X$};
    \node at (2.5,-0.6) {\small $\vx-\vp=s\vu+t\vv$};
  \end{scope}
\end{tikzpicture}
\end{center}

\begin{center}
\renewcommand{\arraystretch}{1.6}
\resizebox{\linewidth}{!}{%
\begin{tabular}{c|cccc}
\toprule
 & \textbf{شکل نرمال} & \textbf{شکل عمومی} & \textbf{شکل برداری} & \textbf{شکل پارامتری}\\
\midrule
\textbf{خط} &
$\begin{cases}\vn_1\cdot\vx=\vn_1\cdot\vp_1\\ \vn_2\cdot\vx=\vn_2\cdot\vp_2\end{cases}$ &
$\begin{cases}a_1x+b_1y+c_1z=d_1\\ a_2x+b_2y+c_2z=d_2\end{cases}$ &
$\vx=\vp+t\vd$ &
$\begin{cases}x=p_1+td_1\\ y=p_2+td_2\\ z=p_3+td_3\end{cases}$\\
\midrule
\textbf{صفحه} &
$\vn\cdot\vx=\vn\cdot\vp$ &
$ax+by+cz=d$ &
$\vx=\vp+s\vu+t\vv$ &
$\begin{cases}x=p_1+su_1+tv_1\\ y=p_2+su_2+tv_2\\ z=p_3+su_3+tv_3\end{cases}$\\
\bottomrule
\end{tabular}}
\end{center}

\begin{intuition}
یک معادله‌ی خطی در $\R^3$ یک صفحه است و یک خط در $\R^3$ \emph{محل برخورد دو صفحه} است؛ برای همین شکل عمومی خط در $\R^3$ دو معادله دارد. پس حل دستگاه $m$ معادله با ۳ مجهول یعنی پیدا کردن محل برخورد $m$ صفحه. مجموعه‌ی جواب می‌تواند تهی، یک نقطه، یک خط یا یک صفحه باشد. شکل برداری/پارامتری دقیقاً همان چیزی است که الگوریتم حل دستگاه به ما می‌دهد: متغیرهای آزاد همان پارامترها ($s,t$) هستند.
\end{intuition}

\subsection{جواب عمومی به شکل برداری-پارامتری}
\begin{note}[آیا می‌توان ترتیب ستون‌ها را عوض کرد؟]
بله، به شرط این‌که حساب کنیم هر ستون مال کدام متغیر است. دستگاه
\[
\begin{cases}ax+by+cz=d\\ a'x+b'y+c'z=d'\\ a''x+b''y+c''z=d''\end{cases}
\quad\text{\fa{با}}\quad
\begin{cases}ax+cz+by=d\\ a'x+c'z+b'y=d'\\ a''x+c''z+b''y=d''\end{cases}
\]
یکی است؛ فقط ترتیب نوشتن جمله‌ها فرق دارد. پس جابه‌جا کردن ستون‌های ماتریس ضرایب یعنی \emph{تغییر نام متغیرها}، و بعد از حل باید نام‌ها را به حالت اول برگرداند. اما ستون آخر ماتریس افزوده (ستون $\vb$) را نباید با ستون‌های ضرایب جابه‌جا کرد. دقت کنید که این یک عمل سطری نیست و در الگوریتم استاندارد از آن استفاده نمی‌کنیم.
\end{note}

\begin{example}
همه‌ی جواب‌های $A\vx=\vb$ را به شکل برداری بنویسید:
$A=\begin{bmatrix}3&5&-4\\-3&-2&4\\6&1&-8\end{bmatrix}$ و $\vb=\begin{bmatrix}7\\-1\\-4\end{bmatrix}$.
\end{example}
\begin{solution}
\[
\left[\begin{array}{ccc|c}3&5&-4&7\\-3&-2&4&-1\\6&1&-8&-4\end{array}\right]
\sim
\left[\begin{array}{ccc|c}1&0&-\frac43&-1\\0&1&0&2\\0&0&0&0\end{array}\right].
\]
پس $x_1=-1+\frac43x_3$، $x_2=2$ و $x_3$ آزاد است. با نوشتن $x_3=t$:
\[
\vx=\begin{bmatrix}x_1\\x_2\\x_3\end{bmatrix}=\begin{bmatrix}-1+\frac43t\\2\\t\end{bmatrix}=\underbrace{\begin{bmatrix}-1\\2\\0\end{bmatrix}}_{\vp}+t\underbrace{\begin{bmatrix}4/3\\0\\1\end{bmatrix}}_{\vv}.
\]
از نظر هندسی، مجموعه‌ی جواب خطی است که از نقطه‌ی $\vp$ می‌گذرد و با $\vv$ موازی است. طبق قضیه‌ی \ref{thm:general-solution}، $t\vv$ همه‌ی جواب‌های $A\vx=\zero$ است.
\end{solution}

\begin{example}
همه‌ی جواب‌های $A\vx=\zero$ را بنویسید: $A=\begin{bmatrix}1&0&-8&-7\\0&1&4&3\\0&0&0&0\end{bmatrix}$.
\end{example}
\begin{solution}
$A$ همین حالا RREF است: $x_1=8x_3+7x_4$، $x_2=-4x_3-3x_4$ و $x_3=s$ و $x_4=t$ آزادند:
\[
\vx=s\begin{bmatrix}8\\-4\\1\\0\end{bmatrix}+t\begin{bmatrix}7\\-3\\0\\1\end{bmatrix}.
\]
مجموعه‌ی جواب یک «صفحه» (دوبعدی) گذرنده از مبدأ در $\R^4$ است.
\end{solution}

\begin{example}
همه‌ی جواب‌های $A\vx=\vb$ را بنویسید، اگر
$\aug{A}{\vb}=\left[\begin{array}{cccc|c}1&2&-1&1&3\\1&1&-1&1&1\\1&3&-1&1&5\end{array}\right]$.
\end{example}
\begin{solution}
\[
\left[\begin{array}{cccc|c}1&2&-1&1&3\\1&1&-1&1&1\\1&3&-1&1&5\end{array}\right]
\xrightarrow[R_3\leftarrow R_3-R_1]{R_2\leftarrow R_2-R_1}
\left[\begin{array}{cccc|c}1&2&-1&1&3\\0&-1&0&0&-2\\0&1&0&0&2\end{array}\right]
\]
\[
\xrightarrow[R_2\leftarrow -R_2]{R_3\leftarrow R_3+R_2}
\left[\begin{array}{cccc|c}1&2&-1&1&3\\0&1&0&0&2\\0&0&0&0&0\end{array}\right]
\xrightarrow{R_1\leftarrow R_1-2R_2}
\left[\begin{array}{cccc|c}1&0&-1&1&-1\\0&1&0&0&2\\0&0&0&0&0\end{array}\right].
\]
پس $x_1=-1+x_3-x_4$، $x_2=2$ و $x_3=s$، $x_4=t$ آزادند:
\[
\vx=\begin{bmatrix}-1\\2\\0\\0\end{bmatrix}+s\begin{bmatrix}1\\0\\1\\0\end{bmatrix}+t\begin{bmatrix}-1\\0\\0\\1\end{bmatrix}.
\]
\end{solution}

\section{جمع‌بندی و تمرین}

\begin{summary}
\begin{itemize}
  \item هر ماتریس با یک و فقط یک ماتریس RREF هم‌ارز سطری است. جای محورها ذاتی ماتریس است.
  \item عملیات سطری روابط خطی بین ستون‌ها را حفظ می‌کنند: $A\vx=\zero\iff (PA)\vx=\zero$.
  \item دستگاه سازگار است $\iff$ ستون آخر ماتریس افزوده محوری نیست. اگر سازگار باشد: بدون متغیر آزاد یک جواب، با متغیر آزاد بی‌شمار جواب.
  \item $A\vx=\vb$ برای هر $\vb$ جواب دارد $\iff$ $A$ در هر سطر محور دارد.
  \item $A\vx=\zero$ جواب نابدیهی دارد $\iff$ متغیر آزاد دارد؛ به‌ویژه اگر $n>m$.
  \item $A\vx=\vb$ جواب دارد $\iff$ $\vb$ ترکیب خطی ستون‌های $A$ است؛ و جواب عمومی = یک جواب خصوصی + جواب‌های همگن.
  \item اگر حذف بدون جابه‌جایی ممکن باشد، $A=LU$ و درایه‌های زیر قطر $L$ همان ضرب‌کننده‌های حذف‌اند. حل دستگاه: $L\vy=\vb$ سپس $U\vx=\vy$.
\end{itemize}
\end{summary}

\begin{exercise}
فرم RREF ماتریس $\begin{bmatrix}1&2&3\\4&5&6\\7&8&9\end{bmatrix}$ را بیابید. ستون‌های محوری کدام‌اند؟
\end{exercise}

\begin{exercise}
به ازای چه مقادیری از $h$ ماتریس $\left[\begin{array}{cc|c}1&h&4\\3&6&8\end{array}\right]$ ماتریس افزوده‌ی یک دستگاه سازگار است؟
\end{exercise}

\begin{exercise}
تجزیه‌ی $LU$ ماتریس $\begin{bmatrix}2&5\\-4&-7\end{bmatrix}$ را بیابید.
\end{exercise}

\begin{exercise}
مجموعه‌ی جواب معادله‌ی $x_1+2x_2-x_3=4$ (در $\R^3$) را به شکل برداری-پارامتری بنویسید. این مجموعه از نظر هندسی چیست؟
\end{exercise}

\begin{exercise}
درست یا نادرست؟ «هر دستگاه همگن با ۳ معادله و ۴ مجهول جواب نابدیهی دارد.»
\end{exercise}

\begin{exercise}
نشان دهید اگر $A\vx=\vb$ دو جواب متمایز $\vu$ و $\vv$ داشته باشد، $A\vx=\zero$ جواب نابدیهی دارد.
\end{exercise}

\subsection*{پاسخ‌های کوتاه}
\begin{enumerate}
  \item $\begin{bmatrix}1&0&-1\\0&1&2\\0&0&0\end{bmatrix}$؛ ستون‌های محوری ستون‌های ۱ و ۲ هستند.
  \item با $R_2\leftarrow R_2-3R_1$ سطر دوم $\begin{bmatrix}0&6-3h&|&-4\end{bmatrix}$ می‌شود. اگر $h=2$، سطر $\begin{bmatrix}0&0&|&-4\end{bmatrix}$ داریم و ناسازگار است؛ پس دستگاه برای $h\neq2$ سازگار است.
  \item $R_2\leftarrow R_2-(-2)R_1$؛ پس $L=\begin{bmatrix}1&0\\-2&1\end{bmatrix}$ و $U=\begin{bmatrix}2&5\\0&3\end{bmatrix}$.
  \item $x_1=4-2x_2+x_3$؛ پس $\vx=\begin{bmatrix}4\\0\\0\end{bmatrix}+s\begin{bmatrix}-2\\1\\0\end{bmatrix}+t\begin{bmatrix}1\\0\\1\end{bmatrix}$: یک صفحه که از $(4,0,0)$ می‌گذرد.
  \item درست؛ طبق نتیجه‌ی \ref{cor:more-unknowns} چون $4>3$.
  \item $A(\vu-\vv)=\vb-\vb=\zero$ و $\vu-\vv\neq\zero$.
\end{enumerate}
